\lecture{15}
\title{Approximation Algorithms (Part I)}

\begin{document}

\textbf{Notation.} We say that $\mathcal{P}$ denotes (optimization) problems, $P$ denotes instances, and $\textsc{Opt}(P)$ denotes the optimal value of an instance $P$.  We say that $\mathcal{A}$ denotes algorithms, and $\mathcal{A}(P)$ denotes the return value of $\mathcal{A}$ with input $P$.

\textbf{Definition.} An algorithm $\mathcal{A}$ for a minimization problem $\mathcal{P}$ is an ``$\alpha$-approximation'' (for $\alpha \geq 1$) if, for all instances $P$ of $\mathcal{P}$, we have $\mathcal{A}(P) \leq \alpha \cdot \textsc{Opt}(P)$.

\textbf{Definition.} For maximization problems, we have $\mathcal{A}(P) \geq \frac{1}{\alpha} \cdot \textsc{Opt}(P)$ instead.  (Always assume $\alpha \geq 1$.)

\begin{problem}{Maximum Matching}
    Given any graph $G$ (not necessarily bipartite), find a maximum matching.
\end{problem}

\begin{remark}
    It turns out this is actually in $\textbf{P}$.  But the point is that the $\textbf{P}$ algorithm is really complicated, whereas\dots
\end{remark}

\begin{solution}{Maximum Matching}
    Read the edges in some order; greedily pick edges.  This is a $2$-approximation.
\end{solution}

\begin{proof}
    The matching found by this algorithm is \emph{maximal}: no edges can be added.
    \begin{quote}
        \textbf{Claim.} Maximal matchings are $2$-approximations to maximum matchings.

        \emph{Proof:} Consider the maximum matching $M^*$ and our maximal matching $M$.

        No edge from $M^*$ can be added to $M$, so for every $e^* \in M^*$, there is some $e \in M$ that disqualifies it.

        However, each $e \in M$ can disqualify at most two $e^* \in M^*$.  So $|M^*| \leq 2 \cdot |M|$. ~ $\square$
    \end{quote}
    The end.  Some questions: Why is the constant $2$ tight?  When do we achieve equality? ~ $\blacksquare$
\end{proof}

\begin{problem}{Vertex Cover}
    Given a graph $G$, find a minimum vertex cover.
\end{problem}

\begin{solution}{Vertex Cover}
    Use the approximation-maximum-matching algorithm to construct a maximal matching.  The vertex cover is the set of all vertices used in a maximal matching.  This is a $2$-approximation.
\end{solution}

\begin{remark}
    When the matching is good, the vertex cover is bad, and vice versa. (!!)
\end{remark}

\begin{proof}
    We claim endpoints of maximal matchings are $2$-approximations to minimum vertex covers.
    
    Why are endpoints of maximal matchings vertex covers? Because.
        
    And the bound is easy: for any matching $M$, we have $|S^*| \geq |M|$, but $|S| = 2 \cdot |M|$. ~ $\blacksquare$
\end{proof}

Assuming widely-believed conjectures (such as $\textbf{P} \neq \textbf{NP}$), this is the best constant we can achieve.

\begin{solution}{Vertex Cover, Badly}
    Greedily pick vertices of maximum degree.  It turns out this is an $O(\log |V|)$-approximation algorithm---pretty bad.  (Note: $O(\log |V|) = O(\log |E|)$ because $|E| = O(|V|^2)$.)
\end{solution}

\begin{proof}
     Say we start with graph $G_0$, and at iteration $i$, we remove the vertex $v_i \in G_{i - 1}$ with degree $d_i$, yielding a graph $G_i$ with $m_i$ edges.  And say $S^*$ is a minimum vertex cover.

     Notice $d_i \geq \frac{m_{i - 1}}{|S^*|}$, because edges in $G_{i - 1}$ map to vertices in $S^*$ so that each vertex in $S^*$ gets at most $d_i$ edges.

     Therefore, $m_{\ell} \leq \left(1 - \frac{1}{|S^*|}\right)^{\ell} \cdot |E|$, because removing vertex $v_i$ removes at least $\frac{1}{|S^*|}$ of all edges in $G_{i - 1}$.
     
     Thus, when $\ell \approx |S^*| \cdot \log |E|$ we have $m_{\ell} \leq 1$, so our vertex cover has size at most $|S^*| \cdot \log|E|$, done. ~ $\blacksquare$
\end{proof}

\begin{remark}
    We cannot do better than a $2$-approximation for vertex cover.

    As a (hard!) exercise: show that the $\alpha = \log |V|$ bound on our vertex-cover approximation algorithm is tight.
\end{remark}

\end{document}